在完成化压缩口径下，越界根对应的“径向偏离量”可被写成一个完全显式的 \(\mathrm{arcosh}\) 闭式；该闭式把离线竖线的水平偏移从叙事性描述压缩为由 \(P_L\) 根位置刚性决定的数值不变量。\footnote{据内部文稿 \cite{Internal2025ResolutionFoldingPhiPi}（2025）在“相位坐标门—有限窗审计—orbit--Levi 刚性”并置且不增设额外公理的口径下，该偏离量与常数层规范归一化一起形成可独立审计的封闭链。}

\begin{corollary}[单尺度完成化的径向偏离量闭式：越界根的离线距离由 \texorpdfstring{$\mathrm{arcosh}$}{arcosh} 刚性给出]\label{cor:xi-single-scale-radial-deviation-arcosh}
\leanverified{paper\_xi\_single\_scale\_radial\_deviation\_arcosh}
在命题 \ref{prop:xi-rh-interval-criterion} 的完成化表示下，令
\[
\Xi_{\Delta,L}(s)=P_L\!\bigl(S_L(s)\bigr),\qquad S_L(s)=L^{s-\frac12}+L^{\frac12-s},
\]
其中 \(P_L\in\RR[S]\)。
取 \(y:=L^{s-\frac12}\)，则 \(S_L(s)=y+y^{-1}\)。
令 \(S_\star\in\RR\) 为 \(P_L\) 的一根（按重数计），并令 \(y_\star^{\pm}\) 为二次方程 \(y+y^{-1}=S_\star\) 的两解：
\[
y_\star^{\pm}=\frac{S_\star\pm\sqrt{S_\star^2-4}}{2}.
\]
则由复对数多值性，对每个符号 \(\pm\) 与 \(k\in\ZZ\) 有一条零点塔
\[
s=\frac12+\frac{\log y_\star^{\pm}+2\pi i k}{\log L},
\]
并且：
\begin{enumerate}
  \item 若 \(S_\star\in[-2,2]\)，则 \(|y_\star^\pm|=1\)，从而该根所诱导的全部零点满足 \(\Re(s)=\tfrac12\)；
  \item 若 \(|S_\star|>2\)，则该根所诱导的离线零点满足
  \[
  \bigl|\Re(s)-\tfrac12\bigr|
  =\frac{\mathrm{arcosh}\bigl(|S_\star|/2\bigr)}{\log L}.
  \]
\end{enumerate}
\end{corollary}

\begin{proof}
由 \(S_L(s)=y+y^{-1}\) 得二次方程 \(y^2-S_\star y+1=0\)，从而 \(y=y_\star^\pm\)。
再由 \(y=L^{s-\frac12}\) 得
\(
\log y=(s-\frac12)\log L-2\pi i k
\)
（\(k\in\ZZ\)），即得到零点塔表达。
\par
当 \(S_\star\in[-2,2]\) 时可写 \(S_\star=2\cos\theta\)，则 \(y_\star^\pm=e^{\pm i\theta}\)，故 \(|y_\star^\pm|=1\)。
当 \(|S_\star|>2\) 时令 \(|S_\star|=2\cosh\eta\)（\(\eta>0\)，即 \(\eta=\mathrm{arcosh}(|S_\star|/2)\)），则两根满足 \(|y_\star^\pm|=e^{\pm\eta}\)。
由 \(|y|=L^{\Re(s)-1/2}\)（命题 \ref{prop:xi-rh-interval-criterion} 的证明）得
\(
|\Re(s)-\tfrac12|
=\eta/\log L
\)，即结论。
\end{proof}
